%  ============================================================================
% CHAPTER 3: RESULTS
% ============================================================================

\section{Results}
Different scenarios were studied and in the end a comparison between different models entropy was done with the same initial conditions.
\subsection{Basic Money Distribution}

In this scenario, we consider a simple model of money distribution among agents without debt. 
The total money of the system is conserved, and no agent can have a negative balance.
The agents engage with each other randomly but always with a fixed amount of money in each transaction.
It is possible to see the results in Figure \ref{fig:NoDebtImage}, where the histogram of the money distribution is plotted along with the theoretical Boltzmann-Gibbs distribution.

\begin{figure}[H]
    \centering
    \includegraphics[width=0.4\textwidth]{NoDebtImage.png}
    \caption{The histogram shows the money distribution of agents after a $10^6$ number of transactions, while the red line represents the theoretical Boltzmann-Gibbs distribution.
                The simulation was performed with $N=3500$, $M = 100 000$, and $\Delta m = 1$. Agents in this scenario can't have debt.}
    \label{fig:NoDebtImage}
\end{figure}



\subsection{Money Distribution with Debt}

The situation now is similar but in this case the agents can have negative balances up to a maximal threshold debt~($md$).
As shown in Fig.~\ref{fig:DebtImage}, after many interactions, the population seems to follow a Boltzmann-Gibbs distribution.
\begin{figure}[H]
    \centering
    \includegraphics[width=0.4\textwidth]{DebtImage.png}
    \caption{The histogram shows the money distribution of agents after a $10^7$ number of transactions, while the red line represents the theoretical Boltzmann-Gibbs distribution.
                The simulation was performed with $N=3500$, $M = 10 000$, and $\Delta m = 1$. Agents in this scenario can have negative balances and the maximal debt is $md = 50$.}
    \label{fig:DebtImage}
\end{figure}


\subsection{Money Distribution with Proportional money transfers}
Let's now consider a scenario where agents exchange only a fraction $\gamma$ of their money.
Each agent can exchange only a fraction of their money, which is determined by the saving propensity.
The results of the simulation are shown in Fig.~\ref{fig:SavingPropensityImage}, where the histogram of the money distribution is plotted along with the theoretical Gamma distribution.
\begin{figure}[H]
    \centering
    \includegraphics[width=0.4\textwidth]{SavingPropensity.png}
    \caption{The histogram shows the money distribution of agents after $10^6$ number of transactions, while the green line represents the theoretical Gamma distribution and the red line represents the Boltzmann-Gibbs distribution with such temperature.
                The simulation was performed with $N=3500$, $M = 1 000 000$, and $\gamma = 1/3$. Agents in this scenario can exchange only a fraction of their money, determined by the saving propensity $\gamma$.}
    \label{fig:SavingPropensityImage}
\end{figure}

